\documentclass[a4paper,12pt]{article}
\usepackage[utf8]{inputenc}
\usepackage[T1]{fontenc}
\usepackage[italian]{babel}
\usepackage{pgfplots}
\usepackage{amsmath}
\usepackage{float}
\usepackage{geometry}
\geometry{margin=2cm}

\pgfplotsset{compat=1.18}

\title{Report 4.1 -- Combinazioni di k e TIMES}

\begin{document}

\maketitle

\newpage

\section{Grafici dei risultati per $k$ e TIMES}

I grafici mostrano rispettivamente l'andamento di:
\begin{itemize}
    \item Avg Profit GAP
    \item Avg Weight GAP
    \item Avg fQ
    \item fQ Min Found
    \item Profit Found
    \item Weight Found
    \item Profit GAP
    \item Weight GAP
    \item n of fQ Min Found
\end{itemize}
al variare di TIMES per diversi valori di $k \in \{1,10,100\}$, confrontando
i risultati usando QALS e non utilizzando QALS.

I risultati riportati sono stati ottenuti utilizzando:

\begin{itemize}
    \item $\lambda = \frac{\sum_{i=1}^{n} p_i}{650 \cdot C}$
    \item l'istanza \texttt{ksp\_1.txt}
\end{itemize}

\begin{figure}[H]
\centering
\begin{tikzpicture}
\begin{axis}[
    width=0.95\textwidth,
    height=0.48\textwidth,
    xmode=log,
    log basis x=10,
    xlabel={TIMES},
    ylabel={Avg Profit GAP},
    title={Avg Profit GAP},
    grid=both,
    legend style={at={(0.5,-0.25)}, anchor=north, legend columns=3},
    xtick={1,10,100,1000},
    xticklabels={1,10,100,1000},
    tick label style={font=\small},
    label style={font=\small},
    title style={font=\small},
]
\addplot[blue, thick, mark=*] coordinates {(1, 101) (10, -15.9) (100, 62.2) (1000, 64.6)};
\addplot[red, thick, mark=square*] coordinates {(1, 92) (10, 83.9) (100, 70.9) (1000, 67.5)};
\addplot[green!60!black, thick, mark=triangle*] coordinates {(1, 128) (10, 60.3) (100, 54.2) (1000, 67.3)};
\addplot[blue, thick, dashed, mark=o] coordinates {(1, -104) (10, -104) (100, -104) (1000, -104)};
\addplot[red, thick, dashed, mark=square] coordinates {(1, -104) (10, -104) (100, -104) (1000, -104)};
\addplot[green!60!black, thick, dashed, mark=triangle] coordinates {(1, -104) (10, -104) (100, -104) (1000, -104)};
\legend{
    QALS ($k=1$),
    QALS ($k=10$),
    QALS ($k=100$),
    NO QALS ($k=1$),
    NO QALS ($k=10$),
    NO QALS ($k=100$)
}
\end{axis}
\end{tikzpicture}
\caption{Andamento di \textit{Avg Profit GAP} al variare di TIMES per diversi valori di $k$.}
\label{fig:avg-profit-gap}
\end{figure}

\begin{figure}[H]
\centering
\begin{tikzpicture}
\begin{axis}[
    width=0.95\textwidth,
    height=0.48\textwidth,
    xmode=log,
    log basis x=10,
    xlabel={TIMES},
    ylabel={Avg Weight GAP},
    title={Avg Weight GAP},
    grid=both,
    legend style={at={(0.5,-0.25)}, anchor=north, legend columns=3},
    xtick={1,10,100,1000},
    xticklabels={1,10,100,1000},
    tick label style={font=\small},
    label style={font=\small},
    title style={font=\small},
]
\addplot[blue, thick, mark=*] coordinates {(1, 14) (10, -94.3) (100, -33) (1000, -38.8)};
\addplot[red, thick, mark=square*] coordinates {(1, -9) (10, -19.4) (100, -24.9) (1000, -33)};
\addplot[green!60!black, thick, mark=triangle*] coordinates {(1, -44) (10, -36.2) (100, -40.4) (1000, -34.7)};
\addplot[blue, thick, dashed, mark=o] coordinates {(1, -76) (10, -76) (100, -76) (1000, -76)};
\addplot[red, thick, dashed, mark=square] coordinates {(1, -76) (10, -76) (100, -76) (1000, -76)};
\addplot[green!60!black, thick, dashed, mark=triangle] coordinates {(1, -76) (10, -76) (100, -76) (1000, -76)};
\legend{
    QALS ($k=1$),
    QALS ($k=10$),
    QALS ($k=100$),
    NO QALS ($k=1$),
    NO QALS ($k=10$),
    NO QALS ($k=100$)
}
\end{axis}
\end{tikzpicture}
\caption{Andamento di \textit{Avg Weight GAP} al variare di TIMES per diversi valori di $k$.}
\label{fig:avg-weight-gap}
\end{figure}

\begin{figure}[H]
\centering
\begin{tikzpicture}
\begin{axis}[
    width=0.95\textwidth,
    height=0.48\textwidth,
    xmode=log,
    log basis x=10,
    xlabel={TIMES},
    ylabel={Avg fQ},
    title={Avg fQ},
    grid=both,
    legend style={at={(0.5,-0.25)}, anchor=north, legend columns=3},
    xtick={1,10,100,1000},
    xticklabels={1,10,100,1000},
    tick label style={font=\small},
    label style={font=\small},
    title style={font=\small},
]
\addplot[blue, thick, mark=*] coordinates {(1, -162.23) (10, -114.79) (100, -121.94) (1000, -128.06)};
\addplot[red, thick, mark=square*] coordinates {(1, -171.98) (10, -130.75) (100, -129.07) (1000, -130.8)};
\addplot[green!60!black, thick, mark=triangle*] coordinates {(1, -123.88) (10, -117.96) (100, -130.94) (1000, -127.52)};
\addplot[blue, thick, dashed, mark=o] coordinates {(1, -330.85) (10, -330.85) (100, -330.85) (1000, -330.85)};
\addplot[red, thick, dashed, mark=square] coordinates {(1, -330.85) (10, -330.85) (100, -330.85) (1000, -330.85)};
\addplot[green!60!black, thick, dashed, mark=triangle] coordinates {(1, -330.85) (10, -330.85) (100, -330.85) (1000, -330.85)};
\legend{
    QALS ($k=1$),
    QALS ($k=10$),
    QALS ($k=100$),
    NO QALS ($k=1$),
    NO QALS ($k=10$),
    NO QALS ($k=100$)
}
\end{axis}
\end{tikzpicture}
\caption{Andamento di \textit{Avg fQ} al variare di TIMES per diversi valori di $k$.}
\label{fig:avg-fq}
\end{figure}

\begin{figure}[H]
\centering
\begin{tikzpicture}
\begin{axis}[
    width=0.95\textwidth,
    height=0.48\textwidth,
    xmode=log,
    log basis x=10,
    xlabel={TIMES},
    ylabel={fQ Min Found},
    title={fQ Min Found},
    grid=both,
    legend style={at={(0.5,-0.25)}, anchor=north, legend columns=3},
    xtick={1,10,100,1000},
    xticklabels={1,10,100,1000},
    tick label style={font=\small},
    label style={font=\small},
    title style={font=\small},
]
\addplot[blue, thick, mark=*] coordinates {(1, -162.23) (10, -229.94) (100, -239.39) (1000, -264.5)};
\addplot[red, thick, mark=square*] coordinates {(1, -171.98) (10, -264.5) (100, -264.5) (1000, -284.04)};
\addplot[green!60!black, thick, mark=triangle*] coordinates {(1, -123.88) (10, -181.75) (100, -249.55) (1000, -249.55)};
\addplot[blue, thick, dashed, mark=o] coordinates {(1, -330.85) (10, -330.85) (100, -330.85) (1000, -330.85)};
\addplot[red, thick, dashed, mark=square] coordinates {(1, -330.85) (10, -330.85) (100, -330.85) (1000, -330.85)};
\addplot[green!60!black, thick, dashed, mark=triangle] coordinates {(1, -330.85) (10, -330.85) (100, -330.85) (1000, -330.85)};
\legend{
    QALS ($k=1$),
    QALS ($k=10$),
    QALS ($k=100$),
    NO QALS ($k=1$),
    NO QALS ($k=10$),
    NO QALS ($k=100$)
}
\end{axis}
\end{tikzpicture}
\caption{Andamento di \textit{fQ Min Found} al variare di TIMES per diversi valori di $k$.}
\label{fig:fq-min-found}
\end{figure}

\begin{figure}[H]
\centering
\begin{tikzpicture}
\begin{axis}[
    width=0.95\textwidth,
    height=0.48\textwidth,
    xmode=log,
    log basis x=10,
    xlabel={TIMES},
    ylabel={Profit Found},
    title={Profit Found: Profitto soluzione con fq minima trovata},
    grid=both,
    legend style={at={(0.5,-0.25)}, anchor=north, legend columns=3},
    xtick={1,10,100,1000},
    xticklabels={1,10,100,1000},
    tick label style={font=\small},
    label style={font=\small},
    title style={font=\small},
]
\addplot[blue, thick, mark=*] coordinates {(1, 97) (10, 370) (100, 178) (1000, 198)};
\addplot[red, thick, mark=square*] coordinates {(1, 106) (10, 198) (100, 198) (1000, 219)};
\addplot[green!60!black, thick, mark=triangle*] coordinates {(1, 70) (10, 119) (100, 194) (1000, 194)};
\addplot[blue, thick, dashed, mark=o] coordinates {(1, 302) (10, 302) (100, 302) (1000, 302)};
\addplot[red, thick, dashed, mark=square] coordinates {(1, 302) (10, 302) (100, 302) (1000, 302)};
\addplot[green!60!black, thick, dashed, mark=triangle] coordinates {(1, 302) (10, 302) (100, 302) (1000, 302)};
\legend{
    QALS ($k=1$),
    QALS ($k=10$),
    QALS ($k=100$),
    NO QALS ($k=1$),
    NO QALS ($k=10$),
    NO QALS ($k=100$)
}
\end{axis}
\end{tikzpicture}
\caption{Andamento di \textit{Profit Found} al variare di TIMES per diversi valori di $k$.}
\label{fig:profit-found}
\end{figure}

\begin{figure}[H]
\centering
\begin{tikzpicture}
\begin{axis}[
    width=0.95\textwidth,
    height=0.48\textwidth,
    xmode=log,
    log basis x=10,
    xlabel={TIMES},
    ylabel={Weight Found},
    title={Weight Found: Peso soluzione con fq minima trovata},
    grid=both,
    legend style={at={(0.5,-0.25)}, anchor=north, legend columns=3},
    xtick={1,10,100,1000},
    xticklabels={1,10,100,1000},
    tick label style={font=\small},
    label style={font=\small},
    title style={font=\small},
]
\addplot[blue, thick, mark=*] coordinates {(1, 87) (10, 279) (100, 129) (1000, 101)};
\addplot[red, thick, mark=square*] coordinates {(1, 110) (10, 101) (100, 101) (1000, 116)};
\addplot[green!60!black, thick, mark=triangle*] coordinates {(1, 145) (10, 125) (100, 142) (1000, 142)};
\addplot[blue, thick, dashed, mark=o] coordinates {(1, 177) (10, 177) (100, 177) (1000, 177)};
\addplot[red, thick, dashed, mark=square] coordinates {(1, 177) (10, 177) (100, 177) (1000, 177)};
\addplot[green!60!black, thick, dashed, mark=triangle] coordinates {(1, 177) (10, 177) (100, 177) (1000, 177)};
\legend{
    QALS ($k=1$),
    QALS ($k=10$),
    QALS ($k=100$),
    NO QALS ($k=1$),
    NO QALS ($k=10$),
    NO QALS ($k=100$)
}
\end{axis}
\end{tikzpicture}
\caption{Andamento di \textit{Weight Found} al variare di TIMES per diversi valori di $k$.}
\label{fig:weight-found}
\end{figure}

\begin{figure}[H]
\centering
\begin{tikzpicture}
\begin{axis}[
    width=0.95\textwidth,
    height=0.48\textwidth,
    xmode=log,
    log basis x=10,
    xlabel={TIMES},
    ylabel={Profit GAP},
    title={Profit GAP: Differenza tra la soluzione con fq minima trovata e la soluzione di DP},
    grid=both,
    legend style={at={(0.5,-0.25)}, anchor=north, legend columns=3},
    xtick={1,10,100,1000},
    xticklabels={1,10,100,1000},
    tick label style={font=\small},
    label style={font=\small},
    title style={font=\small},
]
\addplot[blue, thick, mark=*] coordinates {(1, 101) (10, -172) (100, 20) (1000, 0)};
\addplot[red, thick, mark=square*] coordinates {(1, 92) (10, 0) (100, 0) (1000, -21)};
\addplot[green!60!black, thick, mark=triangle*] coordinates {(1, 128) (10, 79) (100, 4) (1000, 4)};
\addplot[blue, thick, dashed, mark=o] coordinates {(1, -104) (10, -104) (100, -104) (1000, -104)};
\addplot[red, thick, dashed, mark=square] coordinates {(1, -104) (10, -104) (100, -104) (1000, -104)};
\addplot[green!60!black, thick, dashed, mark=triangle] coordinates {(1, -104) (10, -104) (100, -104) (1000, -104)};
\legend{
    QALS ($k=1$),
    QALS ($k=10$),
    QALS ($k=100$),
    NO QALS ($k=1$),
    NO QALS ($k=10$),
    NO QALS ($k=100$)
}
\end{axis}
\end{tikzpicture}
\caption{Andamento di \textit{Profit GAP} al variare di TIMES per diversi valori di $k$.}
\label{fig:profit-gap}
\end{figure}

\begin{figure}[H]
\centering
\begin{tikzpicture}
\begin{axis}[
    width=0.95\textwidth,
    height=0.48\textwidth,
    xmode=log,
    log basis x=10,
    xlabel={TIMES},
    ylabel={Weight GAP},
    title={Weight GAP: Differenza tra la soluzione con fq minima trovata e la soluzione di DP},
    grid=both,
    legend style={at={(0.5,-0.25)}, anchor=north, legend columns=3},
    xtick={1,10,100,1000},
    xticklabels={1,10,100,1000},
    tick label style={font=\small},
    label style={font=\small},
    title style={font=\small},
]
\addplot[blue, thick, mark=*] coordinates {(1, 14) (10, -178) (100, -28) (1000, 0)};
\addplot[red, thick, mark=square*] coordinates {(1, -9) (10, 0) (100, 0) (1000, -15)};
\addplot[green!60!black, thick, mark=triangle*] coordinates {(1, -44) (10, -24) (100, -41) (1000, -41)};
\addplot[blue, thick, dashed, mark=o] coordinates {(1, -76) (10, -76) (100, -76) (1000, -76)};
\addplot[red, thick, dashed, mark=square] coordinates {(1, -76) (10, -76) (100, -76) (1000, -76)};
\addplot[green!60!black, thick, dashed, mark=triangle] coordinates {(1, -76) (10, -76) (100, -76) (1000, -76)};
\legend{
    QALS ($k=1$),
    QALS ($k=10$),
    QALS ($k=100$),
    NO QALS ($k=1$),
    NO QALS ($k=10$),
    NO QALS ($k=100$)
}
\end{axis}
\end{tikzpicture}
\caption{Andamento di \textit{Weight GAP} al variare di TIMES per diversi valori di $k$.}
\label{fig:weight-gap}
\end{figure}

\begin{figure}[H]
\centering
\begin{tikzpicture}
\begin{axis}[
    width=0.95\textwidth,
    height=0.48\textwidth,
    xmode=log,
    log basis x=10,
    xlabel={TIMES},
    ylabel={n of fQ Min Found},
    title={n of fQ Min Found},
    grid=both,
    legend style={at={(0.5,-0.25)}, anchor=north, legend columns=3},
    xtick={1,10,100,1000},
    xticklabels={1,10,100,1000},
    tick label style={font=\small},
    label style={font=\small},
    title style={font=\small},
]
\addplot[blue, thick, mark=*] coordinates {(1, 1) (10, 1) (100, 1) (1000, 2)};
\addplot[red, thick, mark=square*] coordinates {(1, 1) (10, 1) (100, 1) (1000, 1)};
\addplot[green!60!black, thick, mark=triangle*] coordinates {(1, 1) (10, 1) (100, 1) (1000, 2)};
\addplot[blue, thick, dashed, mark=o] coordinates {(1, 1) (10, 10) (100, 100) (1000, 1000)};
\addplot[red, thick, dashed, mark=square] coordinates {(1, 1) (10, 10) (100, 100) (1000, 1000)};
\addplot[green!60!black, thick, dashed, mark=triangle] coordinates {(1, 1) (10, 10) (100, 100) (1000, 1000)};
\legend{
    QALS ($k=1$),
    QALS ($k=10$),
    QALS ($k=100$),
    NO QALS ($k=1$),
    NO QALS ($k=10$),
    NO QALS ($k=100$)
}
\end{axis}
\end{tikzpicture}
\caption{Andamento di \textit{n of fQ Min Found} al variare di TIMES per diversi valori di $k$.}
\label{fig:n-of-fq-min-found}
\end{figure}

\end{document}